% =====================================================================================
% SUITES DU POINT AVEC HUGO -- etat point par point au 13 aout 2026.
%
% VERSION 2 : ajout d'un VERDICT par ligne et d'une ouverture sur les meilleures trouvailles,
% meme gabarit que la note destinee a Caroline Hillairet.
%
% Le point qui structure ce document : l'essentiel des demandes portait sur la RECONCILIATION
% de l'ecart de capital, et la reponse tient dans un residu transforme en table complete.
% =====================================================================================
\documentclass[11pt,a4paper]{article}

\usepackage[utf8]{inputenc}
\usepackage[T1]{fontenc}
\usepackage{lmodern}
\usepackage[french]{babel}
\usepackage[margin=1.7cm]{geometry}
\usepackage{microtype}
\usepackage{array}
\usepackage{booktabs}
\usepackage{longtable}
\usepackage[table]{xcolor}
\usepackage[most]{tcolorbox}
\usepackage{fancyhdr}
\usepackage{lastpage}
\usepackage{eurosym}
\usepackage{enumitem}

\definecolor{navy}{HTML}{1F3864}
\definecolor{lightblue}{HTML}{D6E0F0}
\definecolor{vert}{HTML}{D9EAD3}
\definecolor{ambre}{HTML}{FCE5CD}
\definecolor{gris}{HTML}{EFEFEF}
\definecolor{rouge}{HTML}{F4CCCC}
\definecolor{ordouble}{HTML}{FFF2CC}

\newtcolorbox{cle}[1][]{colback=lightblue!40,colframe=navy,boxrule=0.6pt,
  arc=2pt,left=8pt,right=8pt,top=6pt,bottom=6pt,#1}
\newtcolorbox{phare}[1][]{colback=ordouble,colframe=navy!80,boxrule=0.9pt,
  arc=2pt,left=9pt,right=9pt,top=7pt,bottom=7pt,#1}

\newcommand{\ok}{\cellcolor{vert}\footnotesize\textbf{INTÉGRÉ}}
\newcommand{\autrement}{\cellcolor{vert}\footnotesize\textbf{INTÉGRÉ}\newline\footnotesize\textbf{AUTREMENT}}
\newcommand{\ecarte}{\cellcolor{ambre}\footnotesize\textbf{MESURÉ PUIS}\newline\footnotesize\textbf{ÉCARTÉ}}
\newcommand{\ouvert}{\cellcolor{ambre}\footnotesize\textbf{CHIFFRÉ,}\newline\footnotesize\textbf{À TRANCHER}}
\newcommand{\non}{\cellcolor{rouge}\footnotesize\textbf{NON TRAITÉ}}
\newcommand{\toi}{\cellcolor{gris}\footnotesize\textbf{HORS}\newline\footnotesize\textbf{PÉRIMÈTRE}}
\newcommand{\verdict}[1]{\par\smallskip\emph{\textbf{Verdict.} #1}}

\pagestyle{fancy}
\fancyhf{}
\lfoot{\footnotesize Suites du point avec Hugo --- 13 août 2026}
\rfoot{\footnotesize page \thepage/\pageref{LastPage}}
\renewcommand{\headrulewidth}{0pt}

\title{\vspace{-1.3cm}\color{navy}\bfseries Suites du point tuteur avec Hugo\\[2pt]
  \large Ce qui est intégré, ce qui ne l'est pas, et ce que cela vaut}
\author{Kélian Kaddouri}
\date{13 août 2026}

\begin{document}
\maketitle
\thispagestyle{fancy}
\vspace{-0.7cm}

\begin{phare}
\textbf{Les quatre trouvailles à retenir, si le temps manque.}
\begin{enumerate}[leftmargin=*,itemsep=2pt,topsep=3pt]
\item \textbf{Les cinq milliards manquants ne sont pas une erreur, ils sont le résultat.} La somme
des leviers isolés donne $9\,138$ quand l'écart vaut $14\,139$ : l'écart est l'\emph{interaction},
$+5\,001$~M\euro{} soit $+35\,\%$, et elle est désormais décomposée en quinze termes dont la somme
redonne l'écart à la précision machine. Un résidu est devenu une table.
\item \textbf{C'est la colonne de \emph{fermeture} qu'un plan de remédiation doit citer, pas
l'isolée.} Le même canal vaut $4\,328$~M\euro{} depuis l'état conforme et $8\,775$ depuis l'état
non conforme, parce qu'en le fermant on ferme aussi les croisés qu'il portait. Conséquence de
gestion directe : remédier rapporte davantage à une entité défaillante partout.
\item \textbf{Une paire d'effets croisés est négative, et c'est un résultat.} Propagation
$\times$ accumulation vaut $-330$~M\euro{} sur les seize graines et les deux métriques : les deux
canaux de co-occurrence sont \emph{substituts}, un pilier déjà touché ne pouvant l'être deux fois.
La cascade est super-additive en bloc mais \emph{sous}-additive entre ses deux canaux de
co-occurrence, ce qui borne l'amplitude imputable à la contagion prise en général.
\item \textbf{La borne du retour sur conformité était annoncée à l'envers.} Le mémoire disait
\og plancher \fg{} de la \emph{durée} alors que ce qui est minoré est le \emph{bénéfice} : or
minorer le bénéfice \textbf{majore} la durée. Les $6$ à $35$ ans sont un \textbf{majorant}. Trois
endroits portaient l'erreur, ils sont corrigés.
\end{enumerate}
\end{phare}

\vspace{0.2cm}
\begin{cle}
\textbf{Le fil de ce point.} Quatre demandes sur cinq tournaient autour d'un même constat : la
somme des leviers pris isolément donnait $9$~milliards quand l'écart total en vaut $14$. Chaque
ligne du tableau se termine par un \emph{verdict} qui dit ce que le point vaut, y compris quand la
réponse est \og peu d'enjeu \fg.
\end{cle}

\vspace{0.3cm}
{\footnotesize
\renewcommand{\arraystretch}{1.25}
\begin{longtable}{@{}p{3.4cm} >{\raggedright\arraybackslash\hyphenpenalty=10000\exhyphenpenalty=10000}p{2.4cm} p{8.7cm} p{2.0cm}@{}}
\toprule
\textbf{Demande} & \textbf{Statut} & \textbf{Ce qui a été fait, et ce que cela vaut} &
\textbf{Où le voir} \\
\midrule
\endfirsthead
\toprule
\textbf{Demande \emph{(suite)}} & \textbf{Statut} & \textbf{Ce qui a été fait} &
\textbf{Où} \\
\midrule
\endhead

Retirer les conclusions tirées du bootstrap de $\xi$ &
\ok &
Retirées. L'intervalle $[0{,}30\,;0{,}83]$ \textbf{reste affiché comme un fait} ; ce qui disparaît
est la lecture qu'on en faisait. Deux trouvailles en faisant ce retrait. La slide portait une
valeur \textbf{fausse}, $[0{,}32\,;0{,}84]$ : le $0{,}32$ est la borne basse de la
\emph{delta-méthode} ($[0{,}320\,;0{,}871]$) et le $0{,}84$ n'appartient à aucun des deux
intervalles calculés. Et le même travers vivait dans le mémoire, qui écrivait que $\hat\xi$
\og se stabilise autour de $0{,}6$ \fg{} alors que le balayage le fait descendre à $0{,}38$.
\verdict{La demande était juste et son exécution a rapporté plus qu'elle ne demandait : deux
valeurs invérifiables sont tombées. Le script imprimait d'ailleurs en clair l'interdiction
d'écrire cette phrase, et personne ne l'avait lue.} &
ch.~socle \\

Tableau complet des leviers, et réconciliation avec les $14$~Md\euro &
\ok &
Les seize configurations du treillis sont énumérées et la réconciliation est \textbf{exacte à la
précision machine}. Trois lectures d'un même canal en sortent : isolé, fermeture, Shapley, sommant
respectivement à $9\,138$, $19\,141$ et $14\,139$.
\verdict{\textbf{La meilleure ligne de la liste.} Elle transforme un résidu en table complète, et
elle produit une conséquence de gestion qui n'était pas demandée : c'est la colonne de
\emph{fermeture} qu'un plan de remédiation doit citer, parce qu'un canal fermé ferme aussi les
croisés qu'il portait.} &
ch.~résultats \\

Où passent les $5$~Md\euro{} manquants &
\ok &
Décomposition de Möbius, quinze termes : ordre~1 $9\,138$, ordre~2 $5\,345$, ordre~3 $-691$,
ordre~4 $+346$. \textbf{Mais la réconciliation est une identité, pas une mesure} : ce qu'elle
atteste est que la table est \emph{complète}, pas que chaque terme est précis. Sur seize graines,
neuf termes sur quinze ont un signe résolu, et \textbf{aucun} des cinq d'ordre 3 et 4.
\verdict{Fort, et la réserve compte autant que le résultat. La quasi-annulation des ordres
supérieurs ne se lit \emph{pas} comme cinq mesures ; annoncer l'inverse serait surinterpréter une
identité algébrique.} &
ch.~résultats \\

Effets croisés entre leviers &
\autrement &
Les six paires sont publiées avec leur bruit, et \textbf{il ne faut pas nommer de paire
dominante} : les deux premières se tiennent à $73$~M\euro{} pour des écarts-types de $708$ et
$354$, et leur ordre s'inverse selon la métrique. Ce qui \emph{est} résolu : les trois paires
porteuses passent toutes par la fréquence et font $86\,\%$ de l'ordre~2. Et \textbf{une seule
paire est négative}, propagation $\times$ accumulation.
\verdict{Fort, sur le résultat négatif plus que sur le classement. Deux canaux de co-occurrence
\emph{substituts} bornent ce qu'on peut imputer à la contagion prise en général : c'est le genre
de résultat qu'un jury retient, parce qu'il va contre le sens du modèle.} &
ch.~résultats \\

Clarifier \og portage \fg{} et \og plancher \fg &
\autrement &
Les deux mots étaient mal employés, et la relecture a fait tomber une erreur de sens : le mémoire
écrivait que le \emph{retour} était un plancher alors que ce qui est minoré est le
\emph{bénéfice}. Les trois endroits sont corrigés. Le portage vaut $848$~M\euro/an à $6\,\%$, la
sinistralité évitée $3\,247$~M\euro/an, soit \textbf{$3{,}8$ fois} ce qui est retenu.
\verdict{Fort, et instructif sur la méthode : l'incompréhension en séance ne venait pas d'un
défaut de pédagogie mais d'une \emph{erreur réelle}. Le chiffrage de la perte évitée, affirmée
\og majoritaire \fg{} depuis le début sans jamais être calculée, en est le second gain.} &
ch.~cas d'usage \\

Scénarios par taille d'entité &
\autrement &
La renormalisation proportionnelle suppose une élasticité de $1$ à la taille ; \textbf{l'élasticité
mesurée vaut $0{,}204$}, et les deux canaux estimés sont chacun d'un ordre de grandeur sous
l'unité. Le scénario par taille est donc publié comme \emph{encadrement}, et le chiffre d'entité
reste une \textbf{borne supérieure d'ordre de grandeur}.
\verdict{Utile mais inconfortable, et il faut l'assumer : les deux approches se trompent en sens
opposés, donc il n'existe pas de version \og propre \fg{} du scénario par taille aujourd'hui. Le
correctif est identifié, un plafond de sévérité adossé à l'exposition, mais il relève de la
recalibration que le gel interdit.} &
ch.~résultats \\

Ne pas se brider sur l'affichage de l'incertitude &
\ok &
Suivi dans les deux sens. Toute grandeur simulée est publiée \textbf{avec son bruit}, et les six
postures de report sont définies et chiffrées avec le leur. La table dit que \textbf{la posture la
plus prudente est la moins précise}.
\verdict{Peu d'enjeu de calcul, réel enjeu de crédibilité. Publier \og $1\,247 \pm 24$ \fg{} rend
sans objet toute question sur le troisième chiffre ; retirer le $\pm$ pour faire plus net serait
la dégradation exacte contre laquelle ce mémoire se construit.} &
ch.~robustesse \\

% HARNAIS-HORS-SCRIPT: le -64 % est cite DANS LE BUT DE DIRE qu'aucune sortie versionnee ne le
% reproduit sur cette grandeur. C'est le seul cas ou publier un nombre non confirme est correct :
% le non-confirme EST l'information.
Quel écart entre états publier, et sous quelle définition &
\ouvert &
Deux chiffres sont établis et imprimés : un gain de conformité de $-53\,\%$, soit
$8\,861$~M\euro{} de capital en moins, et un facteur $3{,}34$ entre les deux états de référence.
\textbf{Ils ne mesurent pas la même chose}, parce qu'ils ne déplacent pas les mêmes canaux.
\textbf{Un troisième chiffre circule, de l'ordre de $-64\,\%$, et il ne faut pas l'employer en
l'état} : la vérification faite pour cette note ne le retrouve dans aucune sortie versionnée
\emph{sur cette grandeur}, ses occurrences dans le dossier appartenant au dimensionnement de la
couverture, qui est un autre sujet.
\verdict{À trancher, et c'est la ligne la plus urgente du tableau. Tant qu'elle ne l'est pas, le
mémoire peut publier deux chiffres qui semblent se contredire alors qu'ils répondent à deux
questions différentes. Échanger l'un contre l'autre sans trancher remplacerait un chiffre mal
défini par un autre.} &
ch.~résultats \\

Ancrer les valeurs de $g$ sur des sources publiques &
\ouvert &
Devenu \textbf{optionnel} depuis l'invariance : le capital est croissant en $g$, donc toute
correspondance respectant l'\emph{ordre} produit l'écart quand le forfait reste plat. Seule
l'amplitude est un scénario.
\verdict{Peu d'enjeu désormais, et c'est une bonne nouvelle : l'invariance a transformé un
prérequis en amélioration de présentation. À arbitrer selon le temps disponible, sans risque pour
la thèse.} &
ch.~robustesse \\

Clarification écrite sur le portage, à envoyer &
\toi &
Le contenu existe et est vérifié, dans le mémoire et dans la table ci-dessus.
\verdict{Sans enjeu de méthode, mais à ne pas laisser traîner : c'est le point du call qui n'avait
pas été compris.} &
--- \\

Second codage en aveugle &
\toi &
Le kit est prêt : dix récits, grille identique au gabarit de saisie, formulaire à jour au 13 août
avec la condition d'aveuglement désormais écrite. Il attend un second codeur, environ une heure.
\verdict{Enjeu réel : c'est le seul remède propre au biais de narration, et le mémoire publie
aujourd'hui un accord inter-juges qui n'a qu'un codeur. Une heure de quelqu'un d'autre vaut
plusieurs pages d'argumentation.} &
--- \\
\bottomrule
\end{longtable}
}

\vspace{0.15cm}
% RATTACHEMENT AU HARNAIS, meme raison que pour la note destinee a Caroline Hillairet.
{\scriptsize\noindent\textbf{Sources de cette note :} scripts~\texttt{20} et~\texttt{20b}
(interaction entre piliers), \texttt{25} (réaction du capital au score de conformité),
\texttt{43} et~\texttt{68} (les quatre canaux, leur attribution et leurs croisés),
\texttt{46} et~\texttt{51} (postures de report et leur bruit), \texttt{47} (bootstrap de
l'indice de queue), \texttt{50} (portage et sinistralité évitée), \texttt{66} (invariance en
$g$), \texttt{70} (élasticité à la taille).\par}

\vspace{0.2cm}
\begin{cle}
\textbf{Les deux interactions à ne jamais additionner.} Un point de vocabulaire qui a failli
produire une erreur de lecture. Il y a l'interaction entre les quatre \emph{canaux},
$+5\,001$~M\euro{} de signe résolu, et celle entre les cinq \emph{piliers}, $+1\,395$~M\euro{} de
signe \emph{non} résolu entre graines. Ce sont \textbf{deux partitions du même écart}, pas deux
effets qui s'ajoutent. Le titre d'une slide invitait précisément à les additionner : il est
corrigé.
\end{cle}

\end{document}
